\subsection{The original prescribed rational class on the whole actual normal metric}

\begin{lemma}[Whole original point recovered from its true normalized Gauss lift]
\label{local:actual-normalized-full-point-gauss-inverse}
\lean{EnvelopeProperty.ProjSubvariety.normalizedSimultaneousProductPointAt_gauss, EnvelopeProperty.ProjSubvariety.normalizedSimultaneousBaseRestrictionAt_gauss}
\leanok
\uses{local:actual-normalized-full-point-restriction, local:actual-product-restriction-gauss-inverse}
For an original projectively normal variety and the SAME normalized
simultaneous model, the entire original product projection of each true
normalized Gauss lift is the SAME original product Gauss point. Its
actual original analytic restriction is the SAME original point.
\end{lemma}
\begin{proof}
Compose the established WHOLE normalization-point and blowup-point
projection identities for the actual original Gauss lift. Apply the
proved whole original product restriction identity. No recovery solely
from coordinate values is used, and no perfectness assumption is needed
for these point identities themselves.
\end{proof}

\begin{theorem}[True whole prescribed-class metric from original local convexity]
\label{local:actual-prescribed-class-whole-normal-metric}
\lean{EnvelopeProperty.ProjSubvariety.IsLocallyConvexAll.exists_prescribedClass_full_normalizedRatioMetric}
\leanok
\uses{local:actual-pl-whole-normal-ratio-metric, local:actual-full-point-polarized-ratio-weight, local:actual-normalized-full-point-gauss-inverse}
Let $X$ be the ORIGINAL projectively normal variety over any field,
$a\in\mathbb Q_{>0}$, and let the ORIGINAL function $\varphi$ satisfy
$\operatorname{IsLocallyConvexAll}(a,\varphi)$. There are positive integers
$r,d$ with EXACT $d=r a$, an integer $C$, and actual original product
ideals $I,J$ whose true logs are finite on ALL actual product points.
On the SAME actual normalized model they define the SAME actual
polynomial $O(d)$ tensor Cartier ratio line, with the following behavior.
At EVERY actual original model point $y$, the TRUE canonical ratio weight
plus $C$ is $r\varphi(\operatorname{res}y)$. For EVERY original degree $d$
form $F$ and coordinate-frame index, the TRUE corresponding polarized
section weight is
$r\varphi(\operatorname{res}y)-C-\log|F|_{\operatorname{res}y}$,
including vanishing forms. On each ORIGINAL admissible affine chart for
any such $F$ nonzero modulo the ORIGINAL defining ideal, the SAME actual
section weight at actual normalized Gauss lifts, plus $C$, is convex PL.
No degree nonnegativity, nefness or positivity certificate is assumed or
claimed in this metric-identification theorem.
\end{theorem}
\begin{proof}
Original local convexity first gives original PL. Apply the whole actual
normal ratio-model construction to $a.\operatorname{den}\,\varphi$.
If its actual positive scale is $s$, take
$r=s\,a.\operatorname{den}$ and $d=s\,a.\operatorname{num}$, using the
positive numerator as a natural number. Positivity of $a$ gives positive
$r,d$, and the EXACT rational numerator-denominator identity proves
$d=ra$. Associativity of the actual real scaling gives the whole actual
ratio-weight formula at every model point. Subtract the finite integer
shift from this TRUE equality, then use the full original polarized
section-weight formula to obtain the stated formula for every original
$F$, with no nonvanishing restriction.

Fix an ORIGINAL admissible degree $d$ form nonzero modulo the ORIGINAL
ideal. Apply the ORIGINAL local-convexity hypothesis to the exact
admissible pair $(r,F)$, using $d=ra$, and multiply its convex PL chart
function by the positive rational integer $r$. At every original chart
point, the preceding whole normalized Gauss inverse recovers the SAME
original point. Its form value is nonzero by the genuine original chart
construction. Substitute the TRUE whole ratio metric equality and TRUE
polarized section formula. The extended logarithm there is the original
finite real logarithm. The actual real identity, after multiplying by
$r$, identifies this SAME genuine section weight plus $C$ with the
scaled ORIGINAL admissible chart function. Thus its convex PL property
follows from original local convexity alone. Original class $a$ is kept
throughout; no unrelated large positive twist replaces it.
\end{proof}
